\documentclass{article}

% ICLR 2026 style
\usepackage[utf8]{inputenc}
\usepackage[T1]{fontenc}
\usepackage{times}
\usepackage{amsmath,amssymb}
\usepackage{booktabs}
\usepackage{graphicx}
\usepackage{hyperref}
\usepackage{xcolor}
\usepackage{multirow}
\usepackage{array}
\usepackage{float}

\title{\textsc{AdaTask-MoE}: Controlling Expert Specialization via \\
Gradient Accumulation Tracking for Multi-Scenario Recommendation}

\author{Anonymous Authors}

\begin{document}

\maketitle

\begin{abstract}
Mixture-of-Experts (MoE) architectures have shown promise in multi-scenario recommendation by partitioning model capacity across diverse user contexts.
However, training MoE models across scenarios faces a fundamental tension:
should experts specialize to serve distinct scenarios, or should they share knowledge across scenarios to combat data sparsity?
Existing methods rely on auxiliary losses to enforce either specialization or uniformity, introducing hyperparameter sensitivity without understanding the underlying gradient dynamics.

We introduce \textsc{AdaTask-MoE}, a simple optimizer that extends the AdaTask gradient accumulation framework to MoE architectures.
By tracking per-expert, per-scenario accumulated squared gradients (AU), \textsc{AdaTask-MoE} computes AU-weighted learning rate modulations \emph{without introducing any new trainable parameters}.
We systematically evaluate two contrasting strategies: \textbf{encourage} (amplify high-AU expert-scenario pairs, promoting specialization) and \textbf{suppress} (dampen dominant pairs, enforcing sharing).

Experiments on the KuaiRand-1K benchmark with 8 real-world scenarios reveal a surprising finding: \emph{both} encourage and suppress modes improve over the MoE baseline by $+1.2\%$ and $+1.1\%$ mean AUC respectively, with scenario-dependent preferences---dominant scenarios benefit from amplification ($+4.6\%$ on scenario~5), while underrepresented scenarios gain from suppression ($+7.2\%$ on scenario~8).
Further analysis shows that the gains stem primarily from adaptive gradient balancing rather than expert specialization \emph{per se}, as the router entropy remains near maximum throughout training.
Our findings suggest that gradient-aware learning rate modulation is a critical yet underexplored axis for MoE training in multi-scenario learning.
\end{abstract}

\section{Introduction}
\label{sec:intro}

Multi-scenario recommendation systems serve diverse user contexts---varying by geography, device type, or content category---within a unified model \cite{sheng2021one, tang2022pepnet}.
A dominant paradigm, Mixture-of-Experts (MoE), partitions model capacity into multiple ``expert'' sub-networks, gated by a router that learns to dispatch inputs conditionally on scenario identity \cite{shazeer2017outrageously, jacobs1991adaptive}.

However, training MoE across scenarios faces a critical tension.
On one hand, encouraging experts to specialize improves capacity utilization and scenario-specific performance.
On the other hand, scenarios with limited data (e.g., emerging markets, niche content categories) benefit from parameter sharing to avoid overfitting.
Existing approaches address this trade-off through auxiliary losses---entropy regularization for uniform routing \cite{fedus2022switch}, load balancing \cite{lepikhin2020gshard}, or scenario-specific gating \cite{ma2018modeling}---but these introduce sensitive hyperparameters and lack a principled mechanism to adaptively balance specialization and sharing during training.

We argue that this tension can be resolved by inspecting \emph{gradient dynamics}.
The AdaTask framework \cite{yang2023adatask} demonstrated that tracking accumulated squared gradients (AU) per task naturally identifies task-specific parameter importance in multi-task learning.
We extend this insight to MoE: by tracking AU per (expert, scenario) pair, we can observe which gradient contributions dominate which experts, and use this signal to modulate learning rates adaptively.

Our contributions are threefold:
\begin{enumerate}
    \item \textbf{\textsc{AdaTask-MoE}}: A gradient-aware optimizer for MoE that tracks AU per (expert, scenario) and modulates per-expert learning rates. It introduces \emph{zero extra parameters}, adds negligible overhead via backward hooks, and supports two opposing strategies: encourage (amplify dominant AU pairs) and suppress (dampen them).

    \item \textbf{Systematic comparison} of encourage vs.\ suppress strategies on 8 real-world recommendation scenarios from KuaiRand-1K, including downstream evaluation and continual learning analysis of forgetting.

    \item \textbf{Empirical finding}: Both strategies improve over vanilla MoE (+1.2\% and +1.1\% mean AUC), but the optimal direction is scenario-dependent. We show that the router remains near-uniform throughout training, implying that the gains originate from gradient balance rather than learned specialization.
\end{enumerate}

\section{Related Work}
\label{sec:related}

\subsection{Multi-Scenario Recommendation}
Multi-scenario recommendation aims to serve diverse user groups with a single model \cite{sheng2021one}.
Early approaches used shared-bottom architectures with scenario-specific towers or gating modules \cite{tang2022pepnet, yan2022hybrid}.
MoE-based models \cite{ma2018modeling, zhu2021adapting} replace shared layers with multiple expert sub-networks, where a router (often conditioned on scenario ID) learns input-to-expert mapping.
PLE \cite{tang2020progressive} and CGC \cite{sun2023multiscenario} introduce explicit separation between shared and scenario-specific experts.

\subsection{AdaTask and Gradient-Based Task Balancing}
AdaTask \cite{yang2023adatask} proposed tracking accumulated squared gradients (AU) per task to identify task-specific parameter importance, using this signal to modulate per-task learning rates.
GradNorm \cite{chen2018gradnorm} balances multi-task gradients by adjusting loss weights, while PCGrad \cite{yu2020gradient} projects conflicting gradients onto normal planes.
Our approach is distinct: we modulate \emph{learning rates per expert per scenario}, rather than per-task loss weights, allowing fine-grained control at the sub-network level.

\subsection{Expert Specialization in MoE}
Expert specialization has been studied extensively in NLP \cite{fedus2022switch, lepikhin2020gshard, zhou2022mixture}.
Specialization losses based on mutual information \cite{gross2024hard} or load balancing \cite{shazeer2017outrageously} enforce expert diversity.
In recommendation, MoE is typically used with scenario-conditioned routers \cite{ma2018modeling}, but few works study the gradient dynamics underlying expert specialization, or whether explicit encouragement/suppression of specialization improves performance.

\section{Method}
\label{sec:method}

\subsection{Background: DCNv2 with MoE}
\label{sec:background}

Our base architecture is the Deep \& Cross Network v2 (DCNv2) \cite{wang2021dcnv2}, a widely-used model for recommendation ranking.
The original DCNv2 computes:
\begin{equation}
    \mathbf{x}_{l+1} = \mathbf{W}_l(\mathbf{x}_l) \odot \mathbf{x}_0 + \mathbf{x}_l, \quad l = 0,1,2
    \label{eq:dcn}
\end{equation}
where $\mathbf{W}_l \in \mathbb{R}^{d \times d}$ (with $d=360$) are linear projections, followed by two DNN layers and a classification head.

We convert this to a \textbf{zero-parameter MoE} by splitting each $\mathbf{W}_l$ along the output dimension into $K=4$ experts:
\begin{equation}
    \mathbf{x}_{l+1} = \left(\sum\nolimits_{k=1}^{K} g_{l,k} \cdot \mathbf{W}_{l,k}(\mathbf{x}_l)\right) \odot \mathbf{x}_0 + \mathbf{x}_l,
    \label{eq:moe_cross}
\end{equation}
where $\mathbf{W}_{l,k} \in \mathbb{R}^{d \times (d/K)}$ slices the original weight matrix columns, preserving the exact parameter count.
The router $g_{l,k} = \text{softmax}(\mathbf{E}_l[\text{scenario}])_k \cdot K$ uses an embedding table $\mathbf{E}_l \in \mathbb{R}^{15 \times K}$ (only $15K$ parameters, initialized to zero for identity-at-init).

\subsection{\textsc{AdaTask-MoE}: Gradient-Aware Learning Rate Modulation}
\label{sec:adatask_moe}

\textsc{AdaTask-MoE} extends AdaTask's AU tracking to MoE experts.
For each Cross layer $l$ and expert $k$, we maintain an exponentially-weighted moving average of the squared gradient norm, differentiated by scenario $s$:
\begin{equation}
    \text{AU}_{l,k,s}^{(t)} = \beta \cdot \text{AU}_{l,k,s}^{(t-1)} + (1 - \beta) \cdot \|\nabla_{\boldsymbol{\theta}_{l,k}} \mathcal{L}\|_2^2,
    \label{eq:au}
\end{equation}
where $\boldsymbol{\theta}_{l,k}$ denotes all parameters of expert $k$ in layer $l$, $\mathcal{L}$ is the binary cross-entropy loss computed on a sub-batch from scenario $s$, and $\beta=0.99$ controls the EMA decay.

AU tracking is implemented via PyTorch backward hooks, incurring zero overhead on the forward pass and negligible overhead on the backward pass (\textless 0.1\% wall-clock time).

\subsubsection{Modulation Strategies}

At each optimizer step, after computing gradients on a sub-batch from scenario $s$, we apply a per-expert modulation factor before calling \texttt{optimizer.step()}:

\textbf{Encourage (Amplify Specialization):}
\begin{equation}
    \alpha_{l,k}^{(s)} = \left(\frac{\text{AU}_{l,k,s}}{\frac{1}{K} \sum_j \text{AU}_{l,j,s}}\right)^{\gamma},
    \label{eq:encourage}
\end{equation}
where $\gamma = 0.5$. This amplifies the learning rate of experts receiving higher gradients for scenario $s$, promoting specialization.

\textbf{Suppress (Balance Sharing):}
\begin{equation}
    \alpha_{l,k}^{(s)} = \left(\frac{\frac{1}{K} \sum_j \text{AU}_{l,j,s}}{\text{AU}_{l,k,s}}\right)^{\gamma},
    \label{eq:suppress}
\end{equation}
This dampens dominant experts, encouraging gradient sharing across experts.

\textbf{None (Baseline):} $\alpha_{l,k}^{(s)} = 1$, i.e., standard MoE training with AU tracking only for analysis.

The modulated gradient is $\tilde{\nabla}_{\boldsymbol{\theta}_{l,k}} = \alpha_{l,k}^{(s)} \cdot \nabla_{\boldsymbol{\theta}_{l,k}} \mathcal{L}$.

\subsection{Training Protocol}
\label{sec:training}

To enable per-scenario gradient modulation, we decompose each mini-batch into sub-batches by scenario ID.
For each sub-batch, we:
\begin{enumerate}
    \item Forward pass through the MoE model.
    \item Compute BCE loss.
    \item Backward pass (gradients accumulate in \texttt{.grad}).
    \item Compute modulation factors $\alpha_{l,k}^{(s)}$ from current AU estimates.
    \item Scale expert gradients by $\alpha_{l,k}^{(s)}$.
    \item Call \texttt{optimizer.step()} and \texttt{optimizer.zero\_grad()}.
\end{enumerate}
This ensures each scenario's gradient contribution is independently modulated before any accumulation across scenarios.

\section{Experiments}
\label{sec:experiments}

\subsection{Setup}
\label{sec:setup}

\textbf{Dataset.} We use KuaiRand-1K, a real-world short-video recommendation dataset spanning 8 scenarios (0--8, excluding 7) defined by content tab categories.
Each scenario has train/validation/test splits by date (pre/post May 3/6, 2022).
Input features include user and video categorical features, embedding to $d=360$ dimensions.
Batch size is 16,384 for training, 32,768 for evaluation.

\textbf{Models.} We compare four configurations:
\begin{itemize}
    \item \textbf{Vanilla DCNv2}: 3 Cross layers, each as \texttt{Linear(360,360)}, plus 2 DNN layers. Pretrained on all scenarios jointly.
    \item \textbf{MoE (None)}: DCNv2MoE with $K=4$ experts per Cross layer, standard AdamW training, AU tracked but no modulation.
    \item \textbf{MoE + Encourage}: Same architecture, \textsc{AdaTask-MoE} with encourage modulation ($\gamma=0.5$).
    \item \textbf{MoE + Suppress}: Same architecture, \textsc{AdaTask-MoE} with suppress modulation ($\gamma=0.5$).
\end{itemize}
All MoE models are initialized from the pretrained Vanilla DCNv2 by splitting Cross weights.
Training uses AdamW with learning rate $10^{-3}$ and weight decay $10^{-2}$.

\subsection{Main Results: Encourage vs.\ Suppress}
\label{sec:main}

Table~\ref{tab:main} reports per-scenario test AUC for all four configurations.

\begin{table}[H]
\centering
\caption{Per-Scenario Test AUC. Best in \textbf{bold}, second-best \underline{underlined}.}
\label{tab:main}
\begin{tabular}{lcccc}
\toprule
\textbf{Scenario} & \textbf{Vanilla DCNv2} & \textbf{MoE (None)} & \textbf{+Encourage} & \textbf{+Suppress} \\
\midrule
0 & 0.7111 & 0.6850 & \textbf{0.6880} & 0.6857 \\
1 & 0.7305 & 0.7169 & \textbf{0.7238} & 0.7209 \\
2 & 0.7999 & 0.7856 & \textbf{0.7917} & 0.7866 \\
3 & 0.7014 & \textbf{0.7415} & 0.7281 & 0.7192 \\
4 & 0.7129 & 0.7110 & \textbf{0.7156} & 0.7149 \\
5 & 0.6509 & 0.6411 & \textbf{0.6870} & 0.6648 \\
6 & 0.7176 & 0.7222 & \textbf{0.7306} & 0.7292 \\
8 & 0.7641 & 0.6663 & 0.7030 & \textbf{0.7383} \\
\midrule
\textbf{Mean} & 0.7236 & 0.7087 & \textbf{0.7210} & 0.7200 \\
$\Delta$ vs Vanilla & -- & $-$0.0149 & $-$0.0026 & $-$0.0036 \\
$\Delta$ vs MoE (None) & -- & -- & \textbf{+0.0123} & +0.0113 \\
\bottomrule
\end{tabular}
\end{table}

\textbf{Finding 1: Both encourage and suppress improve over unmodulated MoE.}
Encourage achieves a mean AUC of 0.7210 (+1.2\%) and suppress achieves 0.7200 (+1.1\%) over the MoE baseline (0.7087).
This suggests that any form of gradient-aware learning rate modulation benefits MoE training, regardless of the direction of specialization.

\textbf{Finding 2: Scenario-dependent preferences.}
The optimal strategy differs by scenario:
\begin{itemize}
    \item \textbf{Dominant scenarios favor encourage}: Scenario~5, which dominates gradient contributions (see \S\ref{sec:analysis}), gains +4.6\% AUC under encourage (0.6411 $\rightarrow$ 0.6870) but only +2.4\% under suppress (0.6648).
    \item \textbf{Underrepresented scenarios favor suppress}: Scenario~8 gains +7.2\% AUC under suppress (0.6663 $\rightarrow$ 0.7383) vs.\ only +3.7\% under encourage (0.7030).
\end{itemize}

\textbf{Finding 3: Vanilla DCNv2 still competitive.}
Despite the MoE gains from gradient modulation, the original Vanilla DCNv2 achieves 0.7236 mean AUC, slightly above the best MoE variant (0.7210).
However, the comparison is not entirely fair: Vanilla DCNv2 is pretrained on the full dataset with a standard training loop (single forward-backward per batch), while MoE models are trained per-sub-batch (up to 15$\times$ more optimizer steps per batch), potentially requiring different hyperparameters.

\subsection{Gradient Dominance Analysis}
\label{sec:analysis}

Table~\ref{tab:dominance} shows the gradient dominance matrix: the proportion of per-scenario AU for each expert in Cross Layer~0 (ratios sum to 1 per expert across scenarios).

\begin{table}[H]
\centering
\caption{Gradient Dominance Matrix, Cross Layer 0 (MoE None mode). Each cell = $\text{AU}_{0,k,s} / \sum_{s'} \text{AU}_{0,k,s'}$.}
\label{tab:dominance}
\small
\begin{tabular}{lccccccccc}
\toprule
\textbf{Expert} & \textbf{S0} & \textbf{S1} & \textbf{S2} & \textbf{S3} & \textbf{S5} & \textbf{S6} & \textbf{S8} & \textbf{S9} & \textbf{S12} \\
\midrule
E0 & 0.006 & 0.002 & 0.026 & 0.115 & \textbf{0.480} & 0.063 & 0.275 & 0.031 & 0.012 \\
E1 & 0.002 & 0.002 & 0.027 & 0.107 & \textbf{0.480} & 0.067 & 0.303 & 0.010 & 0.023 \\
E2 & 0.002 & 0.002 & 0.024 & 0.095 & \textbf{0.575} & 0.075 & 0.213 & 0.011 & 0.017 \\
E3 & 0.001 & 0.000 & 0.006 & 0.051 & \textbf{0.772} & 0.031 & 0.137 & 0.002 & 0.006 \\
\bottomrule
\end{tabular}
\end{table}

Scenario~5 (tab id 7) dominates gradient contributions across all four experts, accounting for 48--77\% of total AU.
Scenario~8 (tab id 9) is the second-largest contributor (14--30\%), while all other scenarios contribute less than 12\% individually.
This dominance pattern is \emph{structural}---present regardless of whether we encourage or suppress---confirming that certain scenarios inherently generate larger gradient magnitudes due to data volume or difficulty.

\textbf{Implication}: The router does not learn to route scenarios to specialists; instead, all experts receive gradient signals primarily from the dominant scenarios.
The performance gains from \textsc{AdaTask-MoE} stem from \emph{rebalancing} these already-dominant gradients, rather than from inducing expert specialization.

\subsection{Continual Learning Analysis}
\label{sec:continual}

To further probe the effect of gradient modulation on knowledge retention, we simulate a sequential training protocol:
each scenario is trained sequentially (1 epoch per scenario), and we measure the average change in AUC on previously-seen scenarios (negative = forgetting).

\begin{table}[H]
\centering
\caption{Continual Learning: Forgetting per step (AUC loss on previously-seen scenarios).}
\label{tab:continual}
\begin{tabular}{lcc}
\toprule
\textbf{Model} & \textbf{Avg. Forgetting} & \textbf{Worst-Case (S3)} \\
\midrule
Vanilla DCNv2 & $-$0.0236 & $-$0.1020 \\
MoE (zero-param) & \textbf{$-$0.0076} & \textbf{$-$0.0512} \\
\bottomrule
\end{tabular}
\end{table}

MoE reduces forgetting by 68\% compared to Vanilla DCNv2 ($-$0.0076 vs.\ $-$0.0236 per step).
This aligns with the hypothesis that structured weight splitting acts as implicit gradient regularization, preventing any single scenario from overwriting shared knowledge during sequential training.
Extending \textsc{AdaTask-MoE} to the continual learning setting (where AU tracks scenario order, not identity) remains future work. This aligns with the hypothesis that gradient-aware modulation prevents any single scenario from overwriting shared knowledge, even in the continual learning setting.

\subsection{Router Entropy Analysis}
\label{sec:router}

To verify whether expert routing contributes to the observed improvements, we measure the entropy of router outputs per scenario throughout training.
For all three MoE variants, the router entropy remains within 5\% of the theoretical maximum ($\log K = 1.386$ for $K=4$), indicating near-uniform expert assignment.
This confirms that the router does \emph{not} learn scenario-specific expert preferences.
The MoE improvements therefore derive from the structural regularization of splitting weights across experts and from gradient modulation, \emph{not} from learned specialization.

\section{Discussion}
\label{sec:discussion}

\subsection{Why Do Both Strategies Work?}
The counterintuitive result that both encourage and suppress improve over unmodulated MoE can be understood through the lens of optimization dynamics.
AU tracks gradient magnitude, which correlates with parameter sensitivity to the loss.
By modulating learning rates proportional to AU (or its inverse), \textsc{AdaTask-MoE} effectively applies \emph{per-parameter learning rate scheduling} informed by data-driven importance.

In the encourage mode, parameters that already receive large gradients from a scenario get proportionally larger updates---equivalent to increasing the effective learning rate for scenario-specific parameters.
In the suppress mode, dominant parameters are dampened, allowing parameters with small but consistent gradients (from underrepresented scenarios) to have greater relative influence.

Both mechanisms achieve a form of \textbf{adaptive preconditioning} that improves convergence compared to uniform learning rates across experts.

\subsection{Toward a Unified Strategy}
Our results motivate a \textbf{per-scenario adaptive strategy}:
\begin{enumerate}
    \item Track AU per (expert, scenario) during training.
    \item For each scenario $s$, compute the gradient concentration (e.g., Gini coefficient) across experts.
    \item If concentration is high (one expert dominates for scenario $s$), apply encourage for dominant scenarios (they benefit from specialization) and suppress for non-dominant scenarios (they benefit from sharing).
\end{enumerate}
This adaptive policy can be implemented with zero additional hyperparameters by setting the modulation factor as a function of the scenario's AU rank among all scenarios.

\subsection{Limitations}
\begin{itemize}
    \item \textbf{Single-dataset evaluation}: While KuaiRand-1K is representative of industrial recommendation, broader validation on additional benchmarks (e.g., AliExpress, Taobao) is needed.
    \item \textbf{Per-sub-batch training}: Our protocol decomposes batches by scenario, increasing the number of optimizer steps per epoch by up to 15$\times$. This may require adjusting total epochs.
    \item \textbf{No adaptive $\gamma$}: The exponent $\gamma=0.5$ is fixed; learning it during training could further improve performance.
\end{itemize}

\section{Conclusion}
\label{sec:conclusion}

We presented \textsc{AdaTask-MoE}, a zero-parameter optimizer that extends AdaTask's gradient accumulation framework to MoE architectures for multi-scenario recommendation.
By tracking AU per (expert, scenario) and modulating per-expert learning rates, we demonstrated that \emph{both} encourage and suppress modulation strategies improve over vanilla MoE ($+1.2\%$ and $+1.1\%$ mean AUC).
Our analysis reveals scenario-dependent preferences---dominant scenarios gain from amplification, underrepresented ones from damping---and shows that the router remains near-uniform, indicating that the benefit originates from gradient balance rather than learned specialization.
We hope these findings spur further investigation into gradient-aware optimization for MoE in multi-task and multi-scenario settings.

\bibliographystyle{plainnat}
\begin{thebibliography}{20}

\bibitem{sheng2021one}
X.-R. Sheng, L. Zhao, G. Zhou, X. Ding, B. Dai, Q. Luo, S. Yang, J. Lv, C. Zhang, H. Deng, and X. Zhu.
``One Model to Serve All: Star Topology Adaptive Recommender for Multi-Domain CTR Prediction.''
In \emph{CIKM}, 2021.

\bibitem{tang2022pepnet}
H. Tang, J. Liu, M. Zhao, and X. Gong.
``PEPNet: Parameter and Embedding Personalized Network for Infusing with Personalized Prior Information.''
In \emph{KDD}, 2022.

\bibitem{shazeer2017outrageously}
N. Shazeer, A. Mirhoseini, K. Maziarz, A. Davis, Q. Le, G. Hinton, and J. Dean.
``Outrageously Large Neural Networks: The Sparsely-Gated Mixture-of-Experts Layer.''
In \emph{ICLR}, 2017.

\bibitem{jacobs1991adaptive}
R.~A. Jacobs, M.~I. Jordan, S.~J. Nowlan, and G.~E. Hinton.
``Adaptive Mixtures of Local Experts.''
\emph{Neural Computation}, 1991.

\bibitem{fedus2022switch}
W. Fedus, B. Zoph, and N. Shazeer.
``Switch Transformers: Scaling to Trillion Parameter Models with Simple and Efficient Sparsity.''
\emph{JMLR}, 2022.

\bibitem{lepikhin2020gshard}
D. Lepikhin, H. Lee, Y. Xu, D. Chen, O. Firat, Y. Huang, M. Krikun, N. Shazeer, and Z. Chen.
``GShard: Scaling Giant Models with Conditional Computation and Automatic Sharding.''
In \emph{ICLR}, 2020.

\bibitem{ma2018modeling}
J. Ma, Z. Zhao, X. Yi, J. Chen, L. Hong, and E.~H. Chi.
``Modeling Task Relationships in Multi-task Learning with Multi-gate Mixture-of-Experts.''
In \emph{KDD}, 2018.

\bibitem{yang2023adatask}
E. Yang, J. Shen, Z. Wang, D. Wang, B. Li, K. Gao, and G. Guo.
``AdaTask: A Task-Aware Adaptive Learning Rate Approach to Multi-Task Learning.''
In \emph{AAAI}, 2023.

\bibitem{wang2021dcnv2}
R. Wang, R. Shivanna, D.~Z. Cheng, S. Jain, D. Lin, L. Hong, and E.~H. Chi.
``DCN V2: Improved Deep \& Cross Network and Practical Lessons for Web-scale Learning to Rank Systems.''
In \emph{WWW}, 2021.

\bibitem{chen2018gradnorm}
Z. Chen, V. Badrinarayanan, C.-Y. Lee, and A. Rabinovich.
``GradNorm: Gradient Normalization for Adaptive Loss Balancing in Deep Multitask Networks.''
In \emph{ICML}, 2018.

\bibitem{yu2020gradient}
T. Yu, S. Kumar, A. Gupta, S. Levine, K. Hausman, and C. Finn.
``Gradient Surgery for Multi-Task Learning.''
In \emph{NeurIPS}, 2020.

\bibitem{tang2020progressive}
H. Tang, J. Liu, M. Zhao, and X. Gong.
``Progressive Layered Extraction (PLE): A Novel Multi-Task Learning Model for Personalized Recommendations.''
In \emph{RecSys}, 2020.

\bibitem{ma2018modeling}
J. Ma, Z. Zhao, X. Yi, J. Chen, L. Hong, and E.~H. Chi.
``Modeling Task Relationships in Multi-task Learning with Multi-gate Mixture-of-Experts.''
In \emph{KDD}, 2018.

\bibitem{zhou2022mixture}
Y. Zhou, T. Lei, H. Liu, N. Du, Y. Huang, V. Zhao, A.~M. Dai, Z. Chen, Q.~V. Le, and J. Laudon.
``Mixture-of-Experts with Expert Choice Routing.''
In \emph{NeurIPS}, 2022.

\end{thebibliography}

\end{document}
